%% -*- coding: utf-8 -*-
\documentclass[12pt,pagesize,paper=192mm:108mm,landscape]{scrbook} 
%1920x1080 1280x720
\areaset[current]{192mm}{108mm}
\usepackage{calc}
\usepackage[T2A]{fontenc}
\usepackage[utf8]{inputenc}
\usepackage[english,russian]{babel}
\usepackage{microtype}
\usepackage{misccorr}
\usepackage{cmap}
%\usepackage[unicode=true]{hyperref}
\usepackage{graphicx}
\usepackage{amssymb}
\usepackage{amsmath}
%\usepackage{srcltx}
\usepackage{textcomp}
\usepackage{xspace}
%научные символы и смайлики \smiley \frownie
\usepackage{wasysym}
\usepackage{ccicons}
\begin{document}
\begin{titlepage}
  \vspace*{-0.5em}
  \begin{center}    
    \hspace*{3em}
    \begin{minipage}[t]{1.5em}
      \includegraphics[width=\textwidth]{../BINP-logo}
    \end{minipage}\hfill
    \begin{minipage}{0.115\linewidth}
    \includegraphics[width=\textwidth]{../NSU-logo}
    \end{minipage}
    \hfill
    \hspace*{4.5em}
    \smallskip

    Кафедра теоретической физики физического факультета НГУ

    \Large
    Профессор Грабовский А.\,В.

    \huge
    \textbf{Общая теория относительности}

    \Large
    Лекция № 11
    \vfill

    \small
    \begin{minipage}{0.95\linewidth}
      Наблюдение кластера Пули в оптическом и рентгеновском диапазоне,
      сопоставление наблюдений с реконструкцией распределения массы по
      данным гравитационного линзирования. Интерпретация результата
      как наличие темной материи. Космологический
      принцип. Сопутствующая система координат. Метрика
      Робертсона"=Уокера. Масштабный коэффициент. Собственное
      расстояние на гиперповерхности постоянного времени в метрике
      Робертсона"=Уокера. Скорость разбегания галактик "--- закон
      Хаббла. Постоянная Хаббла через масштабный коэффициент. Связь
      расстояния до галактики, излучившей фотон, с временем излучения,
      красным смещением галактики и масштабным фактором во время
      излучения. Уравнения Фридмана. Уравнение состояния вещества для
      пыли, электромагнитного излучения и темной энергии. Замкнутая,
      плоская и открытая вселенная в зависимости от значения
      плотности.
     \end{minipage}
    \vfill

    \normalsize \ccbysa\hspace{0.5em}  Новосибирск 2022
  \end{center}
\end{titlepage}
\end{document}
